\section{Sequential Deletion and Scaling Protocol}
\subsection{Stateful deletion process}
Sequential unlearning is modeled as a trajectory rather than a collection of unrelated
single deletions.  Starting from $M^{(0)}$, request $F_t$ produces approximate state
$M_u^{(t)}=U(M_u^{(t-1)},F_t)$.  Its counterfactual reference population is independently
trained on
\begin{equation}
  D_r^{(t)} = D \setminus \bigcup_{s=1}^{t} F_s.
\end{equation}
The cumulative deletion manifest is immutable at every step, and all previous forgotten
identifiers remain excluded.  Exact models are not initialized from the approximate state.

\subsection{Trajectory endpoints}
At each step the study records utility, probability divergence, calibration, privacy,
representation, local behavior, and cost.  Drift is analyzed both absolutely and relative
to the contemporaneous exact envelope.  Candidate patterns include monotone accumulation,
nonlinear failure, plateau, and recovery, but none is assumed.  The area above the exact
envelope may be reported as a descriptive trajectory summary only after its weighting and
normalization are fixed.

\subsection{Deletion regimes}
The architecture supports random, high-loss, class-specific, rare-example,
representation-clustered, mislabeled, influence-selected, and sequential selections.
Selection is performed from frozen eligible indices and produces a hashed manifest before
unlearning.  High-loss and influence-based policies must calculate their scores from an
explicit model checkpoint and may not use test labels.  Corruption experiments preserve
both clean and corrupted labels so that removal of harmful data is distinguishable from
ordinary IID deletion.

\subsection{Compute--fidelity frontier}
Cost is measured jointly with fidelity rather than presented as an unconditional speedup.
For method $a$, retraining-relative cost is
\begin{equation}
  \rho_a = \frac{T_a}{\operatorname{median}_i T(M_{r_i})},
\end{equation}
with hardware, device, synchronization policy, peak memory, and measurement boundary
recorded.  Pareto comparisons identify methods for which no alternative is both cheaper
and more faithful.  Wall-clock results remain environment-specific and are not extrapolated
across accelerators without measurement.

\subsection{Escalation across datasets and architectures}
Digits remains a deterministic infrastructure test.  Fashion-MNIST provides a low-cost
neural pilot.  CIFAR-10 is intended to test stronger convolutional architectures,
CIFAR-100 increases class granularity, and a tabular dataset tests whether observations
depend on vision-specific representations.  Dataset inclusion is tied to a generalization
question, not to benchmark count.  The same manifests, method interface, raw schema, and
analysis pipeline are retained across CPU, CUDA, ROCm, Apple accelerators, and remote
hardware; only device adapters and environment records vary.
